\chapter{Supplementary tables}
\label{chap:B-tableaux}

This appendix gives an index of every generated table and the list of result files they come from. Each generated table is written by
\file{report/tools/tables.py} from a JSON result file: no number in it is typed by hand.

\section{RL on large language models in the survey}
\label{sec:b-lrl}

\Cref{tab:sondage-lrl}, which sets the 27 papers of the first 2026 draw on large language models
(family X5 \texttt{llm}, \cref{chap:A-grille}) against classical RL from the same year, is set and
discussed in \cref{chap:04-sondage-resultats}. The 18 \texttt{llm} papers read along the way in the
queue of the enlargement are left out of it (\cref{chap:03-sondage-methode}).

\section{Index of the generated tables}
\label{sec:b-index}

\Cref{tab:b-index} gives, for each generated table, the chapter that discusses it and its source file.

\begin{small}
\begin{longtable}{@{}>{\raggedright\arraybackslash}p{0.19\linewidth}>{\raggedright\arraybackslash}p{0.43\linewidth}>{\raggedright\arraybackslash}p{0.3\linewidth}@{}}
\caption{Where to find each generated table.}\label{tab:b-index}\\
\toprule
Table & Content & Source file \\
\midrule
\endfirsthead
\toprule
Table & Content & Source file \\
\midrule
\endhead
\bottomrule
\endfoot
\Cref{tab:sondage-pineau} & Pineau's criteria, 2018 against classical RL in 2026 &
  \file{extension\_stats.json}, \file{results\_final.json} \\
\Cref{tab:sondage-extension} & Our criteria Y1 to Y12, same columns & \file{extension\_stats.json} \\
\Cref{tab:sondage-lrl} & RL on large models against classical RL, 2026 & \file{extension\_stats.json} \\
\Cref{tab:trajectoire} & Share per practice and per year, trend test & \file{trajectory.json} \\
\Cref{tab:accord} & Agreement of the blind double readings & \file{agree\_v2.json} \\
\Cref{tab:robustesse} & Robustness checks of the survey & \file{robustness.json} and variants \\
\Cref{tab:e1} & E1: thirty PPO runs, false positives & \file{e1.json} \\
\Cref{tab:b1} & B1: the $t$ of 9.09, two ways of computing it & \file{b1\_script.json}, \file{b1\_trpo.json} \\
\Cref{tab:e2} & E2: ranking per task, ten and twenty random seeds, equal budget & \file{e2*.json} \\
\Cref{tab:e2-bonferroni} & E2 with Bonferroni intervals & \file{e2\_bonferroni*.json} \\
\Cref{tab:e3} & E3: PPO PyTorch against PPO JAX & \file{e3*.json} \\
\Cref{tab:e4} & E4: the best three out of ten & \file{e4.json}, \file{e4\_r10.json} \\
\Cref{tab:e5} & E5: how many random seeds recover the order & \file{e5*.json} \\
\Cref{tab:e6} & E6: the tuned baseline & \file{e6*.json} \\
\Cref{tab:e7} & E7: score kept under modified physics & \file{e7*.json} \\
\Cref{tab:dose} & Physics modified by degrees & \file{dose*.json} \\
\Cref{tab:b3} & B3: does the measure change the ranking & \file{b3*.json} \\
\Cref{tab:grpo} & GRPO: every run & \file{summary\_5.json}, \file{gains\_5.json}, \file{extra\_5.json} \\
\Cref{tab:env} & Versions of the environments & \file{envs/*.lock} \\
\Cref{tab:calcul} & Compute per program and per task & \file{sim/logs}, \file{sim/results} \\
\Cref{tab:articles-2026} & The 2026 papers read & \file{final\_2026.json} and queue A3 \\
\Cref{tab:articles-2018} & The 2018 papers read & \file{calib\_2018\_v2.json} \\
\end{longtable}
\end{small}

Sixteen more generated tables go with the complementary analyses, each placed in the chapter of its
analysis. \Cref{tab:b-index-ajouts} refers to each and names the result file
it comes from.

\begin{small}
\begin{longtable}{@{}>{\raggedright\arraybackslash}p{0.22\linewidth}>{\raggedright\arraybackslash}p{0.4\linewidth}>{\raggedright\arraybackslash}p{0.3\linewidth}@{}}
\caption{The generated tables of the complementary analyses.}\label{tab:b-index-ajouts}\\
\toprule
Table & Content & Source file \\
\midrule
\endfirsthead
\toprule
Table & Content & Source file \\
\midrule
\endhead
\bottomrule
\endfoot
\multicolumn{3}{@{}l}{\textit{Survey}}\\
\Cref{tab:accord-intervalles} & Agreement of the pooled double readings, Wilson intervals and kappa
  with a bootstrap interval & \file{survey\_extra.json} \\
\Cref{tab:difference-detectable} & Smallest detectable rise and smallest detectable fall per
  criterion & \file{survey\_extra.json}, \file{extension\_stats.json} \\
\Cref{tab:regle-origine} & The rule that also requires Pineau's figures, applied to the data &
  \file{survey\_extra.json} \\
\Cref{tab:q5-versions} & Q5 according to the version read, proceedings or arXiv &
  \file{q5\_versions.json} \\
\Cref{tab:x8} & X8 against X2, sources of the uncertainty per year &
  \file{x8\_stats.json}, \file{extension\_stats.json}, \file{survey\_extra.json} \\
\multicolumn{3}{@{}l}{\textit{Simulations}}\\
\Cref{tab:hyperparametres} & Settings of the five programs as they ran &
  \file{hyperparams.json} \\
\Cref{tab:regle-tranche} & The \enquote{settled} rule on true null hypotheses &
  \file{rule\_calibration.json} \\
\Cref{tab:e1-grille} & E1 over the sixteen pairs & \file{e1\_grid.json} \\
\Cref{tab:e2-strict} & E2, strict count next to the protocol's count &
  \file{e2\_strict*.json} \\
\Cref{tab:e3-holm} & E3 with the Holm correction & \file{e3\_holm.json} \\
\Cref{tab:e3-puissance} & Power of the E3 test by number of random seeds &
  \file{e3\_holm.json} \\
\Cref{tab:e5-hors-echantillon} & E5 checked out of sample & \file{e5\_holdout.json},
  \file{e5.json} \\
\Cref{tab:e6-vitesse} & Speed of the E6 agents and the agent published by the Zoo & \file{e6\_enjoy.json}, \file{e6\_hub.json}, \file{e6\_speed.json} \\
\Cref{tab:e7-test} & The reversal on \task{Hopper}, tested pair by pair &
  \file{e7\_test*.json} \\
\multicolumn{3}{@{}l}{\textit{GRPO strand}}\\
\Cref{tab:grpo-mesures} & Five measures per configuration, percentile and Student intervals & \file{logs\_stats\_5.json} \\
\Cref{tab:grpo-contraste} & True reward against randomly drawn reward &
  \file{logs\_stats\_5.json} \\
\end{longtable}
\end{small}

\section{The result files}
\label{sec:b-fichiers}

Each JSON result file comes from one script, rerun by \texttt{make analysis}
(\cref{sec:make}). The detail of each criterion and each pair, which the tables summarise, is in
these files (\cref{tab:b-fichiers}).

\begin{small}
\begin{longtable}{@{}>{\raggedright\arraybackslash}p{0.3\linewidth}>{\raggedright\arraybackslash}p{0.64\linewidth}@{}}
\caption{The analysis scripts and the content of their JSON file.}\label{tab:b-fichiers}\\
\toprule
Script & Content of the result file \\
\midrule
\endfirsthead
\toprule
Script & Content of the result file \\
\midrule
\endhead
\bottomrule
\endfoot
\multicolumn{2}{@{}l}{\textit{Survey, folder \file{study/}}}\\
\file{analyze.py} & Shares in the first 2026 draw (\file{results\_2026.json}). \\
\file{analyze\_final.py} & Shares for 2026 and for the 2018 calibration, from the raw readings and
  the harmonisation (\file{results\_final.json}). \\
\file{audit\_stats.py} & Statistical audit of the comparison between 2018 and 2026: Newcombe
  differences, Fisher, Holm, reader agreement (\file{audit\_stats.json}). \\
\file{agree\_v2.py} & Agreement between readers on the raw readings, before harmonisation
  (\file{agree\_v2.json}). \\
\file{extension\_stats.py} & Main comparison: classical RL in 2026 against our reading of 2018,
  criterion by criterion, with Wilson, Newcombe, Fisher and Holm (\file{extension\_stats.json}, with
  variants without the neighbouring model and with the flipped decisions). \\
\file{robustness.py} & Robustness cuts and test of the cost thesis (\file{robustness.json}). \\
\file{title\_filter.py}, \file{title\_filter\_compare.py} & The title filter: share of hidden RL
  and comparison of the twenty hidden papers (\file{title\_filter*.json}). \\
\file{trajectory.py} & The trajectory from 2018 to 2026: shares per year, Cochran-Armitage,
  immediate and delayed gaps (\file{trajectory.json}). \\
\file{falsify\_no.py} & The search against the \enquote{no} decisions (\file{falsify\_no.json}). \\
\file{checklist.py} & The checklists declared at NeurIPS 2025 against our readings
  (\file{checklist.json}). \\
\file{survey\_extra.py} & Double readings with intervals, smallest detectable difference, the rule
  that also requires Pineau's figures, readers and year, mechanical lower bound of X8
  (\file{survey\_extra.json}). \\
\file{q5\_versions.py}, \file{q5\_versions\_stats.py} & Q5 reread in the proceedings and in the
  latest arXiv version, repository dates, four readings of the comparison (\file{q5\_versions.json}). \\
\file{x8\_stats.py} & Criterion X8 against X2, sources of the uncertainty, agreement of the double
  readings (\file{x8\_stats.json}). \\
\file{gemma\_stats.py} & The \texttt{gemma3:12b} reader against our readers: agreement and kappa per
  criterion and per year, year $\times$ reader interaction, 2026 minus 2018 differences read by this
  reader and by ours (\file{gemma\_stats.json}; variant where an invalid decision counts
  as \enquote{no}, \file{gemma\_stats\_invalid\_no.json}). \\
\multicolumn{2}{@{}l}{\textit{Simulations, folder \file{sim/analysis/}}}\\
\file{e1.py}, \file{e1\_curve.py} & E1: thirty PPO runs, Welch test on the final scores, then on the
  curves. \\
\file{b1\_trpo.py}, \file{b1\_script.py} & B1: the readings of the 2018 average $t$ on TRPO and PPO,
  then the computation of Henderson's script. \\
\file{e2.py}, \file{e2\_bonferroni.py} & E2: interquartile mean per task, probabilities of
  improvement, changes of ranking, then the same intervals at the simultaneous Bonferroni level. \\
\file{e3.py} & E3: the two implementations of PPO. \\
\file{e4.py} & E4: mean and spread of the ten runs, then of the best three, reversals. \\
\file{e5.py} & E5: share of draws that recover the order, for 3, 5, 10 and 20 random seeds. \\
\file{e6.py}, \file{e6\_diag.py}, \file{e6\_official.py} & E6: default against Zoo, diagnosis of the
  recipe, official policies evaluated with our protocol. \\
\file{e7.py}, \file{e7\_rank.py} & E7: share of score kept under modified physics, then the
  descriptive ranking. \\
\file{dose.py} & The physics modified by degrees. \\
\file{b3.py} & B3: three measures per run and the reversals between measures. \\
\file{e2\_strict.py} & E2, strict count of the pairs that reverse. \\
\file{e1\_grid.py} & E1 over the sixteen pairs. \\
\file{e3\_holm.py} & E3 with the Holm correction, Hodges-Lehmann difference, power. \\
\file{e7\_test.py} & E7, each pair tested in the training physics and the modified physics. \\
\file{e5\_holdout.py} & E5 out of sample. \\
\file{rule\_calibration.py} & False-positive rate of the \enquote{settled} rule. \\
\file{b1\_factor.py} & B1, the ratio between the $t$ of the script and that of the random seeds. \\
\file{hyperparams.py} & The settings of the programs as they ran (\texttt{make hyperparams}). \\
\file{curves.py} & Learning curves: ranking along the steps, final score against best point,
  autocorrelation. \\
\file{effect\_sizes.py} & Effect sizes in score points. \\
\file{e6\_speed.py}, \file{e6\_enjoy.py}, \file{e6\_hub.py} & E6: speed read in the logs, agents
  replayed with our protocol, agent published by the Zoo. \\
\file{e6\_gamma.py} & E6: the Zoo recipe with $\gamma = 0.99$ (\file{e6\_gamma.json}). \\
\file{c6\_determinism.py}, \file{c2.py} & Determinism of SAC and CrossQ, why the two PPOs diverge.
  Written in advance, run by \texttt{make analysis-suite} (\file{c6\_determinism.json},
  \file{c2.json}). \\
\multicolumn{2}{@{}l}{\textit{GRPO strand, folder \file{sim/llm/}}}\\
\file{summary.py} & Pass@1 at steps 0 and 100 of every run, gains, intervals, sensitivity to the
  missing gain, CPU runs alone (\file{summary*.json}, \file{gains*.json}, \file{extra*.json}). \\
\file{logs\_stats.py} & Greedy accuracy, pass@8, truncated share and accuracy of finished answers,
  Student intervals, sign test, contrast between rewards (\file{logs\_stats\_5.json}). \\
\file{c7\_stats.py} & The analysis, written in advance, of a GRPO rerun on ten random seeds:
  no published result. \\
\end{longtable}
\end{small}

The suffixes of the simulation files: \file{\_s20} for twenty random seeds,
\file{\_r10} for the replication of E4 on random seeds 10 to 19, \file{\_equal} and
\file{\_equal\_s20} for the equal budget. Those of the GRPO strand: no
suffix for three random seeds, \file{\_5} for five.
